\chapter*{Introduction: the program and its arc}
\addcontentsline{toc}{chapter}{Introduction: the program and its arc}
\markboth{Introduction}{}

\section*{Where the dives started}

The frontier study that the dives answer summarises the state of marketing measurement in 2026 in one
organising sentence: \emph{experiments are the ground truth, and everything else is calibrated to them}.
The decade's load-bearing negative results---Lewis and Rao's demonstration that purchase variance swamps
plausible advertising effects, the eBay null of Blake, Nosko and Tadelis, and Gordon et al.'s finding across
hundreds of platform RCTs that even double machine learning on thousands of covariates fails to recover
experimental lift---ended the hope that observational statistics could substitute for randomisation. What
replaced it is a triangulation architecture: lift RCTs and geo experiments as causal anchors, Bayesian media
mix models (MMM) as the cross-channel always-on connective tissue with experiments injected as priors or
likelihood terms, and attribution demoted to a fast tactical signal understood to be non-causal. Three forces
shaped that settlement: privacy engineering destroyed user-level joins and pushed measurement back to
aggregates; open-source tooling (Meridian, Robyn, PyMC-Marketing, GeoLift, Trimmed Match) put the frontier in
every analyst's hands; and the delivery algorithms became part of the measurement problem, contaminating the
experiments meant to measure them and closing feedback loops between measurement and spend.

The study's second document catalogued the open problems at the \emph{seams} of that architecture: ten
engineering problems (E1--E10), ten human-understanding problems (H1--H10), eleven mathematical problems
(M1--M11), and three cross-cutting grand challenges---the unified causal measurement system (M2 + M10 + E6),
the human-mechanism panel (H1 + H3 + H5), and the calibrated simulator (E9 + H7). A third document explained
why most companies have not adopted sophisticated MMM: business-logic barriers (unit economics, incentives,
``precision theatre''), mathematical barriers (weak identification, the impossibility of validation, prior
elicitation) and operational ones (data archaeology, geo data, talent, experiment operations).

The deep-dive program took those catalogues as a queue. Phase~A worked the seams where experiments meet
models (M1, M2, M10, E6, M9); Phase~B worked stability, trust and adoption mechanics (M6/M4, the decision
interface, M3, the operational data problem, M5); Phase~C attempted the grand challenges and the frontier's
edges (M8, E9/H7, M7, H1, efficiency under heavy tails); Phase~D began working the program's own backlog,
returning to identification with the tools built in between. Every fourth run extended the most promising
prior dive instead of opening a new one.

\section*{A map of the chapters}

\Cref{tab:map} lists the chapters with the problem each addresses and its principal dependencies. The
chapters are self-contained, but they were not written in isolation: the dives borrowed each other's
constructions, and several later dives exist only because an earlier one left a specific question in the
backlog. A reader who wants the shortest path through the book's central thread should read
\cref{ch:identification} (what a spend path can identify and what it costs), \cref{ch:transport} (how an
experiment enters a model), \cref{ch:voi} (what the information is worth), \cref{ch:deadband} (when to act on
it) and \cref{ch:feedback} (what the acting does to the measuring). The remaining chapters branch off that
spine.

\begin{table}[htbp]
\centering\small
\caption{Chapter map: the frontier-study problem each chapter addresses, its source dives, and the chapters it most depends on or feeds.}
\label{tab:map}
\begin{tabular}{@{}c p{2.55in} c c p{1.55in}@{}}
\toprule
Ch. & Topic & Problem & Dives & Builds on / feeds \\
\midrule
1 & Identified set of MMM parameters; the price of identification; the spectral comb; the price of the prior & M1, BL1--2, 70, 73, 75 & 01, 20, 21, 22 & feeds 2, 3, 7, 13; uses 7, 9 \\
2 & Transport map from lift tests to MMM parameters & M2 & 02 & uses 1; feeds 3, 4, 14 \\
3 & Decision-calibrated uncertainty and EVSI & M10 & 03, 04 & uses 2; feeds 4, 7, 11, 12 \\
4 & Always-on experiment-to-model fusion & E6 & 05 & uses 2, 3; feeds 7, 13 \\
5 & Geo aggregation and privacy noise & M9 & 06 & uses 1, 2 \\
6 & Refresh stability and analyst degrees of freedom & M6/M4 (barriers) & 07 & uses 1; feeds 7, 9 \\
7 & The reallocation deadband & M9(comm.), B5, H8 & 08, 09 & uses 3, 4, 6; feeds 9, 10, 12, 13, 15 \\
8 & Creative estimands under algorithmic delivery & M3 & 10 & feeds 11 \\
9 & The data-archaeology compiler & O1/O3 & 11 & uses 5, 6; feeds 1, 13 \\
10 & Long-run effects under weak surrogacy & M5 & 12 & uses 2, 3; feeds 7, 14 \\
11 & Impossibility of causal attribution; escape designs & M8 & 13, 14 & uses 8, 3; feeds 4, 7 \\
12 & The calibrated simulator & E9/H7 & 15 & uses 3, 7; feeds 15 \\
13 & Inference inside the feedback loop & M7 & 16, 19 & uses 1, 3, 4, 7, 9; feeds 1 \\
14 & The human-mechanism panel & H1 & 17 & uses 1, 2, 10 \\
15 & Lift under heavy tails & M6 (efficiency) & 18 & uses 7, 12 \\
\bottomrule
\end{tabular}
\end{table}

\section*{Recurring ideas across the dives}

Reading the twenty-two reports together, a small number of ideas recur often enough to count as the
program's findings rather than any single dive's. They are stated here so that the reader can watch for them.

\paragraph{The decision-optimal policy destroys the information it needs.}
\cref{ch:identification} shows that the profit-maximising media plan is, to first order, exactly
non-identifying, and that budget-neutral variation cannot identify the total marginal value of media at all
under a common adstock. \cref{ch:feedback} shows that a certainty-equivalent allocation loop without
exogenous exploration stalls at an uninformative spend where the model's own marginal ROAS echoes the
breakeven target exactly. \cref{ch:fusion} finds that belief-driven experiment selection leaves the very
channel that is about to trap unprobed. The same mechanism appears in each: what is best to do given current
beliefs is what teaches least about them. The escapes are also the same across chapters---exogenous probes with
energy off the seasonal harmonics, switch designs, spend floors and coverage constraints, $t^{-1/4}$
exploration---and \cref{ch:identification} prices them in dollars.

\paragraph{Plateaus: design choice is second-order, the statistic is first-order.}
\cref{ch:voi} finds an EVSI power plateau---among well-powered scalar-readout experiments, which channel,
pulse or window to test is decision-irrelevant---and diagnoses it as a ceiling structure. \cref{ch:identification}
finds the same plateau in the frequency domain (the information functional is flat within 10\% across
probe periods from 13 to 182 weeks). \cref{ch:deadband} finds that the \emph{shape} of the multichannel
inaction region is second-order while using the wrong \emph{statistic} (raw estimates, significance rules,
posterior width) costs tens of percent to multiples. The practical reading is that effort spent optimising a
design within its admissible class buys little, while effort spent on the estimand, the projection, the filter
and the count of experiments buys a lot.

\paragraph{Precisely wrong.}
Several dives find a sophisticated procedure that is more confident than a naive one and more biased:
iROAS used as a prior on marginal ROAS (\cref{ch:transport}, IQR seven times tighter, 43\% off); naive
posterior-as-prior refreshing (\cref{ch:refresh}, coverage 0.24); pseudo-inverses on singular Fisher
matrices (\cref{ch:identification}, \BL{82}); an outcome-selected experiment library that makes a simulator
look better than perfect (\cref{ch:simulator}); a model inside a feedback loop whose marginal ROAS reports
the target it was steered to (\cref{ch:feedback}); a difference-in-means lift estimate that at 200,000 users is
worse than not experimenting (\cref{ch:heavytail}). The adoption study called this precision theatre; the
dives supply the mechanisms.

\paragraph{Exact identities as the working method.}
The program's characteristic move is to find an exact decomposition whose terms name the mechanisms:
$\iROAS = m_\infty\,\rho\,C\,Q$ (\cref{ch:transport}); the divergence-bias identity for creative tests
(\cref{ch:algorithmic}); the leakage law for imputed channels (\cref{ch:archaeology}); the surrogate-index
error as bypass tail minus phantom persistence (\cref{ch:surrogacy}); Muth's invariance
$\Var(\Delta\hat z) = q$ that collapses a partially observed control problem to a fully observed one
(\cref{ch:deadband}); exact comb orthogonality through heterogeneous adstock and the ledger identities
through one vector $d$ and one matrix $\tilde I$ (\cref{ch:identification}); the design theorem for
attribution (\cref{ch:attribution}). Almost all were reproduced by the clean-room verifier to machine
precision, which is where the program's confidence in them comes from.

\paragraph{Every assumption has a price, and the price can be computed.}
\cref{ch:heavytail} prices a margin restriction in variance ratios; \cref{ch:identification} prices a shape
prior in excitation dollars and an unknown saturation curve as the harmonic-collision penalty;
\cref{ch:attribution} prices an interaction cap in aliasing and a holdout-only design in $n^4$ precision loss;
\cref{ch:surrogacy} prices a maximum half-life in the width of the identified set; \cref{ch:aggregation}
prices a privacy budget in posterior width. The unifying object is a cost--precision frontier on which a
decision-maker chooses, and several chapters conclude that the decision layer of \cref{ch:deadband} is where
those prices should be paid.

\paragraph{The program corrected itself.}
Each dive's later rounds refuted some of its own earlier claims, and the compendium keeps them under the
\tagRefuted\ label: the predicted 42\% D-optimality forfeit (\cref{ch:voi}), the cadence law as a forecaster
(\cref{ch:fusion}), the closed-form confounding factor and the ``resolution'' account of block length, the
closed-form quadrature angle, the ``700--2400$\times$ worse'' flighting claim that was in fact exact
singularity, and the admissibility constant (\cref{ch:identification}). The reviewer notes add what the dives
did not catch.

\section*{How to read the numbers}

All figures are from simulations under the model stated in each chapter's setup, with the uncertainty the
dive reported; where a dive reported none, the reviewer note says so. Percentages of ``media spend'' or
``revenue'' are in the dives' synthetic units and should be read as orders of magnitude. Where two dives
disagree, both numbers are given and the disagreement is flagged. Where a result depends on a Fisher-matrix
inversion that has not had the \BL{82} audit, the reviewer note names it. The compendium adds no results of
its own.
